\section{Building and checking the artifacts}\label{app:verification}

\subsection{Shared sources and trust boundary}

The accompanying project uses Lean~4 and Mathlib at version 4.34.1
\cite{Lean4,Mathlib2020,Mathlib4}. The mathematical source is in
\lean{lean/TransposeFollowup/}, with the reused native compiler in
\lean{lean/FixedTapeCompiler/}. The local package \lean{lean/upstream/}
points directly to the sibling first paper's sources. Its build outputs are
kept in the adapter; no definitions or proofs from that paper are copied or
modified. The original frozen statement, checker, full file inventory and
digests are checked by \lean{verification/check_upstream.py}.

The external OpenAI source is pinned to commit
\lean{1581d4cbdd69a0faf34b2089614fe2a814316a6a}. An exact 204-module import
closure, with its original namespace and Apache-2.0 license, is installed by
\lean{verification/fetch_openai_metric.py}. The manifest
\lean{verification/openai-metric-source.json} records import edges and every
source digest. The bibliography and attribution ledger identify the imported scalar
projection theorem. The adapter preserves these sources and uses its own
package configuration.

The proofs use only the standard axioms \lean{propext},
\lean{Classical.choice} and \lean{Quot.sound}. The manuscript-specific audit
runs \texttt{\#print axioms} on all explicitly cited formal declarations,
including definitions and imported endpoints. It checks that every requested
declaration has an audit result and that no other axiom occurs. The formal
source guard additionally checks preservation of the audited mathematical
sources and trust settings. This is an ordinary build and transitive axiom audit, not a fresh
empty-environment replay of every imported proof term.

\subsection{Reproduction}

Keep the two paper folders adjacent in the repository. With the pinned Lean
toolchain available through \lean{elan}, run the following commands from the
follow-up folder. The external source installer can use
\texttt{-{}-source-root PATH} for an offline source checkout; its digest checks
are the same.
\begin{verbatim}
python3 verification/fetch_openai_metric.py
python3 verification/check_upstream.py
cd lean
lake exe cache get
cd ..
python3 verification/manuscript/verify.py
python3 verification/manuscript/build_paper.py
\end{verbatim}
The verification command generates the declaration inventory, runs
\texttt{lake build}, builds all cited modules, and runs the manuscript axiom
audit. It also checks the protected sources and retains the commands, logs and
source hashes. A TeX distribution with \lean{latexmk}, BibTeX and
the packages named in \lean{paper/main.tex} builds the manuscript. The final
PDF is \lean{paper/main.pdf}; intermediate TeX files remain in
\lean{paper/build/}.

\subsection{What the verification establishes}

The declaration inventory covers every formal annotation in the manuscript,
including numbered statements, definitions and cited background endpoints. Generation checks module existence;
the Lean audit checks declaration existence and transitive axioms. A separate
semantic review compares the hypotheses, conclusions and quantifier order of
the English statements with the formal source. The machine-readable contracts
and review resolutions are retained in \lean{verification/manuscript/}.

The earlier library audit checked 6,313 declarations across the six theorem
families; the manuscript audit checks those cited here. Source reviews were
performed by AI agents and are not external human peer review.

Compilation is followed by checks for missing references, citations and
formatting overflow, PDF text inspection and rendered-page review. The source ledger
separately records imported theorems, borrowed arguments and literature
comparisons. Neither rendering checks nor kernel checks establish novelty.
